% !TEX program = lualatex
\documentclass[professionalfonts]{beamer}
\newcommand{\slidemode}{1}  % カラー投影モード
\usepackage{beamer_template}
\usepackage{enumerate}

\newcommand{\LectureNum}{3}
\newcommand{\LectureTitle}{微分法則と積分法則}
\newcommand{\CourseName}{応用数学}
\renewcommand{\TermName}{前期}

\begin{document}

\begin{frame}{第3回　微分法則と積分法則}
  \begin{block}{本回の内容}
  微分法則と積分法則
  \end{block}
  \vfill\hfill{\scriptsize \textcolor{gray}{第2章 ラプラス変換　§1}}
\end{frame}

\begin{frame}{定理：原関数の微分法則}
  \begin{block}{}
  \begin{align*}
    L[f'(t)]    &= sF(s) - f(0) \\
    L[f''(t)]   &= s^{2}F(s) - sf(0) - f'(0) \\
    L[f^{(n)}(t)] &= s^n F(s) - s^{n-1}f(0) - s^{n-2}f'(0) - \ldots - f^{(n-1)}(0)
  \end{align*}
  \end{block}
\end{frame}

\begin{frame}{定理：像関数の高次微分}
  \begin{block}{}
  \[
    L[t^n f(t)] = (-1)^n F^{(n)}(s)
  \]
  \end{block}
\end{frame}

\begin{frame}{定理：原関数の積分法則}
  \begin{block}{}
  \[
    L[\int_{0}^t f(\tau)d\tau] = F(s)/s
  \]
  \end{block}
\end{frame}

\begin{frame}{定理：像関数の積分}
  \begin{block}{}
  \[
    L[f(t)/t] = \int_{s}^\infty F(u)du
  \]
  \end{block}
\end{frame}

\begin{frame}{例題9}
  \begin{exampleblock}{問題}
    微分法則を用いて $L[\cos t]$ を求めよ\\[2pt]
  \end{exampleblock}
\end{frame}

\begin{frame}{例題9　解答}
  $f(t) = \sin t$ とおくと $f'(t) = \cos t,\ f(0) = 0$\\[4pt]
  原関数の微分法則 $L[f'(t)] = sF(s) - f(0)$ を適用\\[4pt]
  $L[\cos t] = s\cdot L[\sin t] - 0 = s\cdot 1/(s^{2}+1)$
  \medskip\textbf{答}：$L[\cos t] = s/(s^{2}+1) \ (s > 0)$
\end{frame}

\begin{frame}{例題10}
  \begin{exampleblock}{問題}
    $L[t^n]$ を求めよ（ n は正の整数）\\[2pt]
  \end{exampleblock}
\end{frame}

\begin{frame}{例題10　解答}
  $f(t) = t^n$ とおくと\\[4pt]
  $f^{(n)}(t) = n!, f(0) = f'(0) = \ldots = f^{(n-1)}(0) = 0$\\[4pt]
  微分法則を n 回適用 $: s^n F(s) - 0 = L[n!] = n!/s$
  \medskip\textbf{答}：$L[t^n] = n!/s^{n+1}  (s > 0)$
\end{frame}

\begin{frame}{例題11}
  \begin{exampleblock}{問題}
    $L[\sin(\alpha t)/t]$ を求めよ（ $\alpha$ は正の定数）\\[2pt]
  \end{exampleblock}
\end{frame}

\begin{frame}{例題11　解答}
  像関数の積分 $: L[f(t)/t] = \int_{s}^\infty F(u)du$ を利用\\[4pt]
  $L[\sin(\alpha t)] = \alpha/(u^{2}+\alpha^{2})$ なので\\[4pt]
  $L[\sin(\alpha t)/t] = \int_{s}^\infty \alpha/(u^{2}+\alpha^{2}) du = [\arctan(u/\alpha)]_{s}^\infty$\\[4pt]
  $= \pi/2 - \arctan(s/\alpha)$
  \medskip\textbf{答}：$L[\sin(\alpha t)/t] = \pi/2 - \arctan(s/\alpha) = \arctan(\alpha/s) \ (s > 0)$
\end{frame}

\begin{frame}{演習問題}
  \begin{exampleblock}{自分で解いてみよう}
  \begin{description}
    \item[問13] $L[t \cos(\alpha t)]$ を求めよ
    \item[問14] $f'(t)+2f(t)=t, f(0)=0$ を満たすとき、 $L[f(t)]$ を求めよ
    \item[問15] $L[t^n]$ を求めよ（ n は正の整数）（像関数の高次微分を使用）
    \item[問16] $L[(e^{2t} - e^t)/t]$ を求めよ
  \end{description}
  \end{exampleblock}
\end{frame}

\end{document}
